%%%%%%%%%%%%%%%%%%%% WINTER -- COMPACT COLM-STYLE DRAFT %%%%%%%%%%%%%%%%%%%%%%
%
% ACL style controls margins. Do not add geometry/fullpage packages.
%
%%%%%%%%%%%%%%%%%%%%%%%%%%%%%%%%%%%%%%%%%%%%%%%%%%%%%%%%%%%%%%%%%%%%%%%%%%%%%%%

\documentclass[11pt]{article}

\usepackage[review]{acl}

\usepackage{times}
\usepackage{latexsym}
\usepackage[T1]{fontenc}
\usepackage[utf8]{inputenc}
\usepackage{microtype}
\usepackage{booktabs}
\usepackage{graphicx}
\graphicspath{{figures/}}
\usepackage{multirow}
\usepackage{xcolor}
\usepackage{amsmath}
\usepackage{array}
\usepackage{tabularx}
\usepackage{url}
\usepackage{hyperref}
\usepackage[capitalize,noabbrev]{cleveref}

\definecolor{darkblue}{rgb}{0,0,0.5}
\hypersetup{colorlinks=true, citecolor=darkblue, linkcolor=darkblue, urlcolor=darkblue}

\newcommand{\pass}{\textsc{pass}}
\newcommand{\fail}{\textsc{fail}}
\newcommand{\inconclusive}{\textsc{inconclusive}}
\newcommand{\vibetest}{\textsc{WINTER}}
\newcommand{\qwen}{Qwen3.6}

\title{\vibetest: Written Intent Testing with Evidence\\for Repositories}

\author{Anonymous ACL submission}

\begin{document}
\maketitle

\begin{abstract}
Software correctness remains difficult to validate at the repository level,
especially when intent is specified only in natural language.
At the same time, highly capable LLM agents increasingly modify codebases from
high-level instructions.
We study \vibetest{} (Written Intent Testing with Evidence for Repositories), a
way to test repository-level behavior from natural-language specifications
instead of formal predicates.
In our setting, a testing agent works in a sandboxed copy of a repository,
turns an informal pipeline property into evidence-gathering actions, and
returns \pass/\fail/\inconclusive{}, cited evidence, and a continuous score.
This paper evaluates \vibetest{} on ML repositories, where pipeline intent is
often written informally.
On 1,125 synthetic Kaggle (repository, property) pairs, \vibetest{} with
\qwen{} reaches 0.802--0.837 selective macro F1 and outperforms open-ended
review baselines.
On 1,920 real Kaggle pairs, a conservative sampled fail audit gives
0.931--0.938 macro F1.
The error analysis matters as much as the headline score: among 53 usable
high-confidence synthetic disagreements, 29 are model misses and 24 are
property-operationalization mismatches.
For natural-language repository tests, selective accuracy is only part of the
story. We also have to check how the benchmark labels interpreted each
property.
\end{abstract}

\section{Introduction}

Repository-scale software remains hard to test, even as AI coding agents make
it easier to create large codebases.
Checking code correctness relies on methods for validating that the developer's
intent is reflected in the code, but this remains difficult and costly in
practice~\citep{zeller2009why}.
Even in controlled settings such as secure-development contests, where teams
are explicitly incentivized to build correct, secure systems, submissions
routinely contain vulnerabilities~\citep{ruef2016bibifi}.

At the same time, software engineering is being reshaped by highly capable AI
agents~\citep{metr2025longtasks}.
Benchmarks such as the SWE-bench family~\citep{jimenez2024swebench,
swebench-verified} task models with resolving real-world bug reports in
repositories, and recent analyses find that agentic workflows are increasingly
effective at solving these tasks~\citep{xia2024agentless,guo2025repoaudit}.
Software is already being developed with little or no hand-written code.
Examples include a full tensor library~\citep{vibetensor}, a C compiler
~\citep{carlini2026building}, and internal tools at OpenAI~\citep{
lopopolo2026harness}.
That trend makes the testing problem harder to ignore: more code can be
produced, but intent still has to be checked.

Existing correctness tools cover parts of this problem.
They do not directly handle the setting we study.
Static program analysis is ideally sound with respect to a formal semantics,
but it requires explicit formal specifications, often trades completeness for
tractability~\citep{cousot1979systematic,nielson1999principles}, and claimed
soundness can fail due to simplifying formalization assumptions~\citep{
java-unsound}.
Recent LLM-assisted query systems such as IRIS~\citep{li2024iris},
QLCoder~\citep{wang2025qlcoder}, and KNighter~\citep{yang2025knighter} reduce
manual effort but still depend on formalized properties that may miss
higher-level intent.
Testing and fuzzing explore concrete executions and can uncover defects, but
they take manual work to extend to new properties or
repositories~\citep{manes2018fuzzing}.
The missing case is high-level intent that has not yet been turned into a
formal specification.

We ask: \textit{How can we test whether a repository matches high-level intent
when the intent is written only in natural language?}

We introduce \vibetest{} for this setting.
\vibetest{} operates in a sandboxed copy of the repository, translates an
informal property into analysis and evidence-gathering actions, and returns a
\pass/\fail/\inconclusive{} decision with an evidence report and confidence
score.
Although the method is repository-general, this paper studies it on ML pipeline
properties in Kaggle-style repositories, where high-level requirements such as
no leakage, correct metrics, deterministic evaluation, and consistent device
placement are usually informal.
The results are useful, but they also expose a labeling problem.
\vibetest{} outperforms reviewer-style and direct-property probes; among the
remaining errors, many come from the synthetic labels and the evidence rubric
interpreting the same natural-language property differently.

\begin{figure*}[t]
  \centering
  \includegraphics[width=0.95\linewidth]{vibetest_evidence_verification.png}
  \caption{\vibetest{} takes a repository and natural-language ML-pipeline
  property, gathers evidence, and returns a decision, score, and evidence
  report. Audits check whether predicted failures are supported by the cited
  repository context.}
  \label{fig:overview}
\end{figure*}

\section{Task and Method}
\label{sec:method}

A case is a pair $(R,p)$ of an ML repository and a natural-language property.
The output is $y \in \{\pass,\fail,\inconclusive\}$ plus evidence $E$.
\pass{} means the repository satisfies the property with support; \fail{} means
it violates the property with support; \inconclusive{} means evidence is
insufficient.
This differs from open-ended review, where any bug may be reported, and from
formal analysis, where the property has already been encoded.
In our setting, the central difference is that the property is specified
informally, not as a formal predicate.

\paragraph{Phase 1: Analysis.}
The agent first gets oriented in the ML repository.
It identifies entry points, notebooks, configuration files, data-loading paths,
training code, evaluation code, and artifacts relevant to the property.
The output of this phase is a short plan for where evidence is likely to be
found.

\paragraph{Phase 2: Evidence gathering.}
The agent then inspects code, tests, logs, notebooks, and data-processing logic
to find evidence for or against the property.
When useful, it can run lightweight checks to validate behavior, reproduce a
claim, or expose a violation.
The goal is concrete repository evidence tied to files, outputs, or execution
traces. General code review is outside the task.

\paragraph{Phase 3: Decision.}
Finally, the agent aggregates the evidence and emits a verdict.
\pass{} means the property is supported by evidence, \fail{} means the
repository violates it, and \inconclusive{} means the evidence is too weak or
incomplete.
The agent operates in a sandboxed copy of the repository and may modify files
to probe behavior, but those modifications do not persist to the original
repository.
It also emits \texttt{case\_score} in $[0,1]$, where low values indicate
confidence in \pass{} and high values confidence in \fail{}.
For a two-sided threshold $t$, we accept \pass{} if
\texttt{case\_score} $\leq 1-t$, accept \fail{} if
\texttt{case\_score} $\geq t$, and abstain otherwise.

We test three prompt-example settings: \textbf{ex0} has no examples,
\textbf{ex10} includes ten representative failure examples, and \textbf{ex20}
includes twenty.
The examples are meant to calibrate the evidence standard.
They do not contain answers for the evaluated cases.

\section{Experiments}
\label{sec:experiments}

\paragraph{Setup.}
The synthetic benchmark contains 75 Kaggle-style notebooks across Titanic,
Disaster Tweets, and Diabetic Retinopathy, each evaluated against 15
ML-pipeline properties (1,125 cases).
Labels come from controlled bug injection.
The real benchmark contains 121 public Kaggle notebooks: 50 Titanic, 50
Disaster Tweets, and 21 Diabetic Retinopathy notebooks, evaluated against up to
16 properties (1,920 cases)~\citep{kaggle_titanic,kaggle_disaster_tweets,
kaggle_diabetic_retinopathy}.
Real Kaggle repositories lack ground-truth property labels, so we audit sampled
\fail{} verdicts.

\begin{table}[t]
  \centering
  \small
  \begin{tabular}{llrr}
    \toprule
    Benchmark & Dataset & Props & Repos \\
    \midrule
    \multirow{3}{*}{Synthetic} & Titanic & 15 & 25 \\
     & Disaster Tweets & 15 & 25 \\
     & Diabetic & 15 & 25 \\
    \midrule
    \multirow{3}{*}{Real} & Titanic & $\le$16 & 50 \\
     & Disaster Tweets & $\le$16 & 50 \\
     & Diabetic & $\le$16 & 21 \\
    \bottomrule
  \end{tabular}
  \caption{Synthetic labels come from controlled injection. Real results use
  sampled fail audits.}
  \label{tab:datasets}
\end{table}

\paragraph{Methods.}
We compare \vibetest{} with reviewer-style prompts, a Codex reviewer,
TrainCheck~\citep{traincheck}, and \textbf{Direct Property}.
Reviewer modes perform open-ended review and then map findings to properties.
Direct Property asks an agent to judge one property directly with repository
access but without the full \vibetest{} evidence rubric or examples.
Codex and TrainCheck are execution/format-fragile at scale: in our canonical
run they commit parseable verdicts on 5.5\% and 0\% of cases.
Direct Property, Codex-style review, and \vibetest{} are additionally compared
on the same five-repository Titanic subset.

\paragraph{Research questions.}
We organize the results around four questions.
RQ1 compares \vibetest{} with baseline methods.
RQ2 studies the effect of representative examples in the prompt.
RQ3 audits whether predicted failures are actually supported by evidence, and
uses the audit to compare real and synthetic Kaggle.
RQ4 checks how sensitive the results are to the model backbone and dataset.

\paragraph{Metrics and audits.}
Synthetic metrics use selective macro F1 on covered cases and coverage.
Real macro F1 is computed under a conservative fail audit: sampled \fail{}
decisions are checked by a separate LLM verifier, while \pass{} is taken from
the model output.
We also report worst-class selective error, score distributions, and a sampled
audit of high-confidence synthetic disagreements.

\section{Results}
\label{sec:results}

\paragraph{RQ1: Benchmark and baselines.}
\Cref{tab:headline} and \Cref{fig:synthetic-f1} report that \vibetest{} improves
over full synthetic baselines.
\vibetest{} reaches 0.802 macro F1 for ex0, 0.813 for ex10, and 0.837 for
ex20, while reviewer baselines reach 0.669--0.722.
Codex and TrainCheck are reported at native committed coverage because most
outputs are not parseable committed verdicts.
On the matched five-repository Titanic probe, Direct Property with Qwen Flash
reaches 0.486 failure recall, Codex review reaches 0.135, and \vibetest{} ex10
reaches 0.811 (\Cref{tab:probe}, \Cref{fig:probe-baselines}).
The gain does not come only from asking a stronger model to inspect the
property. The structured evidence procedure changes what the model reports.

\begin{table*}[t]
  \centering
  \scriptsize
  \setlength{\tabcolsep}{3pt}
  \begin{tabular}{llrrrrrcl}
    \toprule
    Benchmark & Method & Ex. & Tests & Cov. & Macro F1 & Out/VF & 95\% CI & Labels \\
    \midrule
    Synthetic & \vibetest{} static & 0 & 1125 & 0.560 & 0.802 & 19.5K & {\scriptsize[0.77,0.83]} & Synthetic GT \\
    Synthetic & \vibetest{} static & 10 & 1125 & 0.643 & 0.813 & 24.5K & {\scriptsize[0.78,0.84]} & Synthetic GT \\
    Synthetic & \vibetest{} static & 20 & 1125 & 0.562 & 0.837 & 21.8K & {\scriptsize[0.81,0.87]} & Synthetic GT \\
    Synthetic & Reviewer 0 & -- & 1125 & 0.973 & 0.669 & 3.8K & {\scriptsize[0.64,0.70]} & Synthetic GT \\
    Synthetic & Reviewer 1 & -- & 1125 & 0.963 & 0.722 & 14.8K & {\scriptsize[0.69,0.75]} & Synthetic GT \\
    Synthetic & Reviewer 2 & -- & 1125 & 0.982 & 0.711 & 23.1K & {\scriptsize[0.68,0.74]} & Synthetic GT \\
    Synthetic & Codex reviewer & -- & 1125 & 0.055 & 0.419 & -- & {\scriptsize[0.33,0.54]} & Synthetic GT \\
    Synthetic & TrainCheck & -- & 1125 & 0.000 & -- & -- & -- & Synthetic GT \\
    Real & \vibetest{} static & 0 & 1920 & 0.822 & 0.931 & -- & -- & Conservative audit \\
    Real & \vibetest{} static & 10 & 1920 & 0.869 & 0.933 & -- & -- & Conservative audit \\
    Real & \vibetest{} static & 20 & 1920 & 0.871 & 0.938 & -- & -- & Conservative audit \\
    \bottomrule
  \end{tabular}
  \caption{Headline results. Synthetic rows report selective macro F1 at each
  method's selected operating point. Real rows use the conservative sampled
  fail audit. Out/VF is output tokens per true predicted failure when measured
  output usage and labels are available.}
  \label{tab:headline}
\end{table*}

\begin{figure}[t]
  \centering
  \includegraphics[width=\linewidth]{paper_synthetic_selective_f1_coverage.pdf}
  \caption{Synthetic selective macro F1 versus coverage. Reviewer
  \inconclusive{} outputs are folded into \pass{} before sweeping a
  pseudo-score, which is charitable to the baselines. Codex and TrainCheck are
  low-coverage native points in \Cref{tab:headline}.}
  \label{fig:synthetic-f1}
\end{figure}

\begin{figure*}[t]
  \centering
  \includegraphics[width=\linewidth]{paper_probe_baseline_metrics.pdf}
  \caption{Matched five-repository probe comparing Direct Property, Codex
  review, and \vibetest{}. \vibetest{} ex10 has the best macro F1 and failure
  recall while keeping high coverage.}
  \label{fig:probe-baselines}
\end{figure*}

\begin{table}[t]
  \centering
  \scriptsize
  \setlength{\tabcolsep}{3pt}
  \begin{tabular}{@{}>{\raggedright\arraybackslash}p{0.36\linewidth}rrrrr@{}}
    \toprule
    Method & Cov. & F1 & P & R & Out/VF \\
    \midrule
    Direct Qwen Flash & 0.840 & 0.620 & 0.750 & 0.486 & 42.6k \\
    Direct GPT-4.1 mini & 0.973 & 0.610 & 0.824 & 0.378 & 8.0k \\
    Codex reviewer & 0.440 & 0.349 & 0.625 & 0.135 & 2.8k \\
    \vibetest{} ex0 & 0.840 & 0.681 & 0.840 & 0.568 & 12.0k \\
    \vibetest{} ex10 & 0.920 & 0.791 & 0.833 & 0.811 & 12.0k \\
    \vibetest{} ex20 & 0.893 & 0.775 & 0.794 & 0.730 & 14.7k \\
    \bottomrule
  \end{tabular}
  \caption{Matched Titanic probe. P/R denote failure precision and recall.
  Out/VF is output tokens per true predicted failure.}
  \label{tab:probe}
\end{table}

\paragraph{RQ2: Example-guided evidence calibration.}
Adding representative examples changes both coverage and failure recall.
On full synthetic Kaggle, ex10 gives the highest coverage among \vibetest{}
settings (0.643), while ex20 gives the highest macro F1 (0.837).
On the matched probe, ex10 also gives the strongest failure recall (0.811).
In this run, the examples seem to set the evidence standard instead of simply
making the model produce longer answers.
\texttt{case\_score} enables selective prediction, but the score is not
perfectly calibrated: worst-class selective error remains nontrivial even after
abstention (\Cref{fig:synthetic-error}), and score ECDFs contain clusters near
confident and ambiguous regions (\Cref{fig:score-hists}).

\begin{figure}[t]
  \centering
  \includegraphics[width=\linewidth]{paper_synthetic_max_error_coverage.pdf}
  \caption{Synthetic worst-class selective error versus coverage on a linear
  scale. Lower is better.}
  \label{fig:synthetic-error}
\end{figure}

\begin{figure*}[t]
  \centering
  \includegraphics[width=\linewidth]{paper_case_score_ecdf.pdf}
  \caption{Empirical case-score distributions. Steep segments indicate score
  clusters; flat segments indicate rarely used score ranges.}
  \label{fig:score-hists}
\end{figure*}

\paragraph{RQ3: Evidence quality and label disagreement.}
An independent verifier confirms 306 of 335 checkable synthetic \fail{}
evidence reports (91.3\%, Wilson CI [0.88,0.94]).
Counting 87 no-checkable-evidence failures against the model lowers fail
precision from 0.79 to 0.73.
On real Kaggle, the conservative fail-audit macro F1 is stable across example
settings: 0.931--0.938 at 0.822--0.871 coverage
(\Cref{tab:headline}, \Cref{fig:real-f1}).
A full-context reaudit accepted 26 of 30 conservative false failures as
supported, so the headline real-data numbers are likely conservative
(\Cref{tab:real-sensitivity}).

Real Kaggle looks stronger than synthetic Kaggle, so we audit high-confidence
synthetic disagreements.
After excluding one parse error, 53 cases remain: 29 are real \qwen{} misses
and 24 are property-operationalization mismatches.
Roughly 45\% of sampled disagreements are therefore not model errors alone. The
injected label and evidence-based judgment use different property boundaries.

\begin{figure}[t]
  \centering
  \includegraphics[width=\linewidth]{paper_real_conservative_f1_coverage.pdf}
  \caption{Conservative macro F1 versus coverage on real Kaggle.}
  \label{fig:real-f1}
\end{figure}

\begin{table}[t]
  \centering
  \small
  \begin{tabular}{lr}
    \toprule
    Examples & Conservative macro F1 \\
    \midrule
    0 & 0.931 \\
    10 & 0.933 \\
    20 & 0.938 \\
    \bottomrule
  \end{tabular}
  \caption{Real Kaggle headline values. Full-context reaudit is treated only
  as an optimistic sensitivity check.}
  \label{tab:real-sensitivity}
\end{table}

\begin{table*}[t]
  \centering
  \small
  \setlength{\tabcolsep}{4pt}
  \begin{tabularx}{\textwidth}{@{}p{2.3cm}p{1.4cm}p{2.25cm}p{2.2cm}X@{}}
    \toprule
    Type & Dataset & Repo & Property & Audit finding \\
    \midrule
    Model miss & Diabetic & \texttt{SYN-D1} & Parameters updated &
    Parameters are frozen before an initial training call; later unfreezing
    does not erase that frozen training stage. \\
    Operationalization mismatch & Diabetic & \texttt{SYN-D1} & Replaceable loops &
    The synthetic label treats benign file/path loops as failures; the evidence
    rubric targets active numerical loops replaceable by vectorization. \\
    Operationalization mismatch & Diabetic & \texttt{SYN-D2} & Train-only augmentation &
    The label says \pass, but the evidence rubric flags stochastic transforms
    reused for train and validation loaders. \\
    \bottomrule
  \end{tabularx}
  \caption{Representative synthetic audit examples.}
  \label{tab:synthetic-examples}
\end{table*}

\paragraph{RQ4: Variation across model and dataset.}
Across datasets, \qwen{} and GPT-5-mini are close at native operating points
(\Cref{tab:models}).
Both are weakest on Diabetic Retinopathy and strongest on Disaster Tweets.
The pattern is simple: task modality and repository structure matter about as
much as model choice.
For that reason, we report dataset-level behavior alongside the aggregate
score.

\begin{table}[t]
  \centering
  \small
  \begin{tabular}{llrrr}
    \toprule
    Dataset & Model & Cov. & F1 & Tok/prop \\
    \midrule
    Titanic & \qwen{} & 0.720 & 0.762 & 2.8k \\
    Titanic & GPT-5-mini & 0.795 & 0.765 & 3.5k \\
    Tweets & \qwen{} & 0.835 & 0.795 & 3.1k \\
    Tweets & GPT-5-mini & 0.731 & 0.756 & 3.1k \\
    Diabetic & \qwen{} & 0.765 & 0.749 & 3.0k \\
    Diabetic & GPT-5-mini & 0.704 & 0.718 & 3.1k \\
    Mean & \qwen{} & 0.773 & 0.769 & 2.9k \\
    Mean & GPT-5-mini & 0.743 & 0.746 & 3.2k \\
    \bottomrule
  \end{tabular}
  \caption{Cross-model comparison at native non-\inconclusive{} coverage.
  GPT-5-mini lacks \texttt{case\_score}, so selective peaks are not compared.}
  \label{tab:models}
\end{table}

\section{Discussion and Limitations}
\label{sec:discussion}

\vibetest{} is not a formal verifier or generic bug finder.
It is a repository-level test for a specified property.
In this paper the properties are ML-pipeline properties, and the useful behavior
comes from forcing the agent to gather evidence for that property instead of
writing a general review.
The reviewer baselines are worse because they are too open-ended.
TrainCheck and Codex expose a different problem: on messy repositories, even
getting a committed, auditable verdict is difficult.
The Direct Property probe is the closest ablation.
It keeps repository access but removes the evidence rubric and examples, and it
loses most clearly on failure recall.

The main limitation is audit scope.
The synthetic disagreement audit is sampled and LLM-adjudicated, not a complete
human relabeling.
The real Kaggle audit samples predicted failures, so it measures sampled fail
precision but does not fully bound missed violations among predicted
\pass{} cases.
Confidence intervals resample cases independently even though properties within
a repository are correlated.
Finally, the same-checkpoint Direct Property baseline is only a five-repository
probe.
A full run would make the baseline comparison stronger.

\section{Related Work}

\paragraph{Repository agents.}
SWE-bench, SWE-bench Verified, Agentless, and RepoAudit evaluate or automate
repository-scale repair and auditing~\citep{jimenez2024swebench,
swebench-verified,xia2024agentless,guo2025repoaudit}.
\vibetest{} instead evaluates a specified natural-language property.

\paragraph{Program analysis and ML debugging.}
Static analysis and CodeQL require formalized properties
~\citep{cousot1979systematic,nielson1999principles,codeql}; LLM-assisted query
generation lowers that burden~\citep{li2024iris,wang2025qlcoder,yang2025knighter}.
ML debuggers such as UMLAUT, TrainCheck, model assertions, and TorchQL target
fixed failure classes~\citep{schoop2021umlaut,traincheck,kang2020assertions,
naik2024torchql}.
\vibetest{} targets properties before they are formalized.

\paragraph{LLM verification.}
Verifiers, process supervision, self-refinement, and LLM judges use separate
reasoning or feedback to improve or evaluate outputs
~\citep{cobbe2021verifier,lightman2024process,madaan2023selfrefine,
weng2023selfverification,zheng2023judge,tan2025judgebench}.
Here, the auditable object is an evidence-backed repository decision.

\section{Conclusion}

\vibetest{} makes repository-level ML intent testable from natural-language
properties. Across synthetic and real Kaggle repositories, structured evidence
gathering improves selective F1 and failure recall over reviewer-style baselines
while producing auditable reports.
The audits add the main lesson: many remaining disagreements come from how a
property is operationalized. Synthetic labels, prompts, and repository evidence
can encode different thresholds for the same violation.
That is a central feature of natural-language testing. \vibetest{} exposes
those boundaries, measures their effect, and turns informal ML correctness
claims into reviewable repository evidence.

\bibliography{vibetest}

\appendix

\section{ML Pipeline Property List}
\label{app:properties}

The property set covers leakage, validation-only model selection, parameter
updates, data load/preprocess code, augmentation, class imbalance,
filename/path features, split definitions, metric definitions, label
randomization, simple baselines, tiny-batch overfit checks, loss curves,
dataloader shuffle policy, deterministic evaluation, device/dtype consistency,
secret use, and loops replaceable by matrix operations.

\section{Real Audit Examples}

\refstepcounter{table}
\label{tab:real-audit-examples}
\begin{center}
  \centering
  \scriptsize
  \setlength{\tabcolsep}{2.5pt}
  \begin{tabularx}{\linewidth}{@{}>{\raggedright\arraybackslash}p{0.18\linewidth}>{\raggedright\arraybackslash}p{0.22\linewidth}>{\raggedright\arraybackslash}p{0.20\linewidth}X@{}}
    \toprule
    Type & Repo & Property & Audit finding \\
    \midrule
    True fail & \texttt{REAL-T1} & Dataloader shuffling &
    Train and validation dataloaders both set \texttt{shuffle=True}. \\
    True fail & \texttt{REAL-T2} & Metric definitions &
    ROC-AUC uses hard class predictions instead of scores. \\
    Context flip & \texttt{REAL-T3} & Replaceable loops &
    Full context contains loop-based \texttt{predict} despite a vectorized
    forward pass. \\
    Context flip & \texttt{REAL-T4} & Device placement &
    Full context contains mixed \texttt{net.to(device)}, \texttt{.cuda()},
    \texttt{.cpu()}, and CPU tensor usage. \\
    \bottomrule
  \end{tabularx}
  \smallskip
  {\footnotesize \textbf{Table~\thetable:} Representative real Kaggle audit
  examples. All rows are Titanic.}
\end{center}

\section{Artifact Map}

Supporting material is in \texttt{supporting\_docs/}. The package also includes
the figure-regeneration and analysis scripts.

\end{document}
